\section*{Chapitre \ref{ch:b1:active-margins} --- Les marges actives}

\begin{solution}{exo:b1:active-margins:1}
(a) $2\pi G\,\Delta\rho\,h = 2\pi\times\num{6.674e-11}\times1700\times3000 =
\qty{2.14e-3}{m.s^{-2}}$~: la pesanteur serait plus basse d'environ
\qty{214}{\milli\gal}. (b) Pour \qty{3.5}{km},
$2\pi\times\num{6.674e-11}\times1700\times3500 = \qty{2.50e-3}{m.s^{-2}}$,
\qty{250}{\milli\gal}~: la taille du minimum observé, donc la fosse est maintenue bien
hors de l'équilibre isostatique. (c) La formule de la plaque suppose que la roche manquante
s'étend sans limite de tous côtés~; une fosse large de dizaines de kilomètres seulement
retire bien moins de masse autour du point de mesure, donc son eau à elle seule
donne moins de \qty{250}{\milli\gal}. Le reste du déficit se trouve plus profond~:
toute la plaque est fléchie vers le bas sous la fosse, et sa \omterm{def:b1:layered-earth:layers}{croûte}, plus légère que
le \omterm{def:b1:layered-earth:layers}{manteau}, prend la place du \omterm{def:b1:layered-earth:layers}{manteau} à sa base.
\end{solution}

\begin{solution}{exo:b1:active-margins:2}
(a) $\qty{82}{mm.\year^{-1}} = \qty{82}{km.\Ma^{-1}}$, donc \qty{820}{km} en
\qty{10}{\Ma}. (b) $820\sin\ang{45} = \qty{580}{km}$. (c) Le chemin jusqu'à
\qty{600}{km} est $600/\sin\ang{45} = \qty{849}{km}$, parcouru en
$849/82 = \qty{10.3}{\Ma}$.
\end{solution}

\begin{solution}{exo:b1:active-margins:3}
$g = 110/\tan\delta$~: \qty{191}{km} pour \ang{30}, \qty{110}{km} pour \ang{45},
\qty{64}{km} pour \ang{60}. L'arc au-dessus du panneau à faible \omterm{def:b1:geological-maps:attitude}{pendage} est le plus éloigné
de sa fosse.
\end{solution}

\begin{solution}{exo:b1:active-margins:4}
(a) Convergence~: $\qty{2.7e6}{m^2.\year^{-1}}/\qty{5.1e7}{m} =
\qty{0.053}{m.\year^{-1}}$, environ \qty{53}{mm.an^{-1}}. Expansion~:
$\qty{3.0e6}{m^2.\year^{-1}}/\qty{6.7e7}{m} = \qty{0.045}{m.\year^{-1}}$, environ
\qty{45}{mm.an^{-1}}. Sur une longueur totale plus courte, les limites de subduction
consomment du plancher océanique plus vite que les dorsales n'en font, mais en consomment moins au total.
(b) Les \qty{0.3}{km^2} restants environ sont repris le long des autres limites
convergentes du modèle, comme les collisions continentales et les zones convergentes
dans les océans, où les plaques sont raccourcies plutôt que subduites~; et
les dorsales ne sont pas les seules limites divergentes~: des rifts continentaux et certaines
transformantes s'ouvrent aussi un peu.
\end{solution}

\begin{solution}{exo:b1:active-margins:5}
(a) $\bar x = 350$, $\bar z = 290$~; les écarts de $x$ sont $-150$, $-50$,
$50$, $150$ et ceux de $z$ sont $-140$, $-40$, $40$, $140$, donc
$S_{xx} = 2(150^2 + 50^2) = \num{50000}$ et
$S_{xz} = 2(150\times140 + 50\times40) = \num{46000}$. Alors $s = 0.92$,
$a = 290 - 0.92\times350 = \qty{-32}{km}$, et le \omterm{def:b1:geological-maps:attitude}{pendage} est
$\arctan 0.92 = \ang{42.6}$. (b) $z = 0$ en $x = 32/0.92 = \qty{34.8}{km}$
de la fosse. (c) Par les points extrêmes la pente est $280/300 = 0.933$, un
\omterm{def:b1:geological-maps:attitude}{pendage} de \ang{43.0}~: proche, mais il n'utilise que deux des quatre séismes.
\end{solution}

\begin{solution}{exo:b1:active-margins:6}
(a) $50/\tan\ang{25} + 60/\tan\ang{45} = 107.2 + 60.0 = \qty{167.2}{km}$.
(b) Un panneau de \ang{45} sur toute la longueur donne \qty{110}{km}~: le segment
superficiel doux élargit la distance de \qty{57}{km}. (c) Le segment profond est
$z = 50 + (x - 107.2)$~; il atteint $z = 0$ en $x = \qty{57.2}{km}$, vers l'intérieur des terres de
la fosse, comme la droite des Kouriles à \qty{53}{km}~: un segment superficiel doux
explique le décalage des Kouriles. (d) $107.2 + 60/\tan\ang{60} = 107.2 + 34.6 =
\qty{141.9}{km}$.
\end{solution}

\begin{solution}{exo:b1:active-margins:7}
(a) Le rayon est la moitié du \omterm{def:b1:geological-maps:attitude}{pendage}, et $r = R\,\theta$ avec $\theta$ en radians~:
\ang{22.5} et $6371\times0.3927 = \qty{2502}{km}$ pour \ang{45}~; \ang{15} et
\qty{1668}{km} pour \ang{30}~; \ang{35} et \qty{3892}{km} pour \ang{70}.
(b) Les arcs au-dessus de panneaux à faible \omterm{def:b1:geological-maps:attitude}{pendage}. (c) La prédiction est
$44.9/2 = \ang{22.45}$~; l'ajustement, \ang{21.9}, est
$0.55/22.45 = \qty{2.4}{\percent}$ plus petit.
\end{solution}

\begin{solution}{exo:b1:active-margins:8}
(a) $q/k$~: $0.035/2.5 = \qty{0.014}{K.m^{-1}} = \qty{14}{K.km^{-1}}$ sous
l'\omterm{def:b1:active-margins:active-margin}{avant-arc}, $0.080/2.5 = \qty{32}{K.km^{-1}}$ sous l'arrière-arc.
(b) $10 + 14\times30 = \qty{430}{\celsius}$ et
$10 + 32\times30 = \qty{970}{\celsius}$. (c) Une partie du flux de chaleur de surface est
produite par la radioactivité dans la \omterm{def:b1:layered-earth:layers}{croûte} elle-même, donc moins de chaleur traverse les
niveaux plus profonds et le gradient diminue avec la profondeur~: une extrapolation linéaire du
gradient de surface surestime la température en profondeur.
\end{solution}

\begin{solution}{exo:b1:active-margins:9}
(a) \qty{32}{mm.\year^{-1}} fait \qty{32}{km.\Ma^{-1}}~; la demi-vitesse est de
\qty{16}{mm.an^{-1}}. (b) $100/32 = \qty{3.1}{\Ma}$. (c) La vitesse n'a pas besoin d'avoir
été constante~; une part de la largeur
peut venir de l'étirement et du rifting de l'arc avant le début de l'expansion~;
et l'axe peut avoir sauté ou s'être ouvert asymétriquement.
\end{solution}

\begin{solution}{exo:b1:active-margins:10}
(a) Le chemin est $600/\sin\ang{45} = \qty{849}{km}$, parcouru à
\qty{82}{km.\Ma^{-1}} en \qty{10.3}{\Ma}, soit \qty{3.27e14}{s}.
(b) $\sqrt{\num{e-6}\times\num{3.27e14}} = \qty{1.8e4}{m}$, environ
\qty{18}{km}. (c) $(\qty{5e4}{m})^2/\qty{e-6}{m^2.s^{-1}} = \qty{2.5e15}{s}$,
environ \qty{79}{\Ma}. (d) La chaleur n'a atteint que les \qty{18}{km} extérieurs de chaque
côté du panneau~: son cœur arrive à \qty{600}{km} presque aussi froid qu'il a quitté
la fosse, plus froid et plus résistant que le \omterm{def:b1:layered-earth:layers}{manteau} qui l'entoure, ce qui explique qu'il puisse
encore se rompre en séismes à cette profondeur.
\end{solution}

\begin{solution}{exo:b1:active-margins:11}
(a) À la distance $L$ derrière le front, la surface se tient $L\tan\alpha$ au-dessus de
l'altitude du front et la base se trouve $L\tan\beta$ en dessous~: l'arrière du
prisme est haut de $L(\tan\alpha + \tan\beta)$, et le triangle a l'aire
$\frac12 L\times L(\tan\alpha + \tan\beta)$. (b) $\tan\ang{4} + \tan\ang{6} =
0.175$. Le prisme reçoit $\qty{1}{km}\times\qty{40}{km.\Ma^{-1}} =
\qty{40}{km^2.\Ma^{-1}}$ de section, donc $A = 40t$ et
$L = \sqrt{2A/0.175}$~: \qty{21.4}{km} après \qty{1}{\Ma}, \qty{42.8}{km}
après \qty{4}{\Ma}, \qty{67.6}{km} après \qty{10}{\Ma}. $L$ croît comme
$\sqrt t$~: quatre fois plus de temps donne deux fois la largeur. (c) $A$ est divisée par deux et $L$
divisé par $\sqrt2$~: \qty{15.1}{km}, \qty{30.2}{km} et \qty{47.8}{km}.
\end{solution}

\begin{solution}{exo:b1:active-margins:12}
(a) $\partial F/\partial a = -2\sum(z_i - a - s x_i) = 0$ donne
$n\bar z = na + sn\bar x$, c'est-à-dire $\bar z = a + s\bar x$. (b) Chaque $z_i$ et
$\bar z$ augmentent de \qty{10}{km}, donc les écarts $z_i - \bar z$ et $S_{xz}$
sont inchangés~: la pente est la même, et $a = \bar z - s\bar x$ augmente de
\qty{10}{km}. (c) Les écarts $x_i - \bar x$ sont inchangés, donc la pente est
la même~; $\bar x$ diminue de 20, donc $a$ augmente de $20s$. (d) Ni l'une ni l'autre~:
le \omterm{def:b1:geological-maps:attitude}{pendage} ne dépend que de la pente. Des erreurs communes à tous les points déplacent la
droite sans la faire tourner.
\end{solution}

\begin{solution}{pb:b1:active-margins:1}
\textbf{1.} Les huit points sont proches d'une droite de \omterm{def:b1:geological-maps:attitude}{pendage} d'environ
\ang{45}, de \qty{96}{km} de profondeur à \qty{155}{km} de la fosse à
\qty{573}{km} de profondeur à \qty{609}{km}.

\textbf{2.} $\sum x_i = 3152$, $\sum z_i = 2751$, $\sum x_i^2 = \num{1442910}$,
$\sum x_i z_i = \num{1290548}$; $\bar x = \qty{394.0}{km}$,
$\bar z = \qty{343.875}{km}$.

\textbf{3.} $S_{xx} = 1442910 - 8\times394^2 = \num{201022}$ et
$S_{xz} = 1290548 - 8\times394\times343.875 = \num{206654}$. Alors
$s = 206654/201022 = 1.028$ et $a = 343.875 - 1.02802\times394 =
\qty{-61.2}{km}$~: $z = -61.2 + 1.028\,x$.

\textbf{4.} $\arctan1.028 = \ang{45.8}$.

\textbf{5.} À $x = 61.2/1.028 = \qty{59.5}{km}$ de la fosse vers l'intérieur des terres. Un panneau
qui aurait quitté la fosse à ce \omterm{def:b1:geological-maps:attitude}{pendage} atteindrait la surface là~; il l'atteint à
la fosse, donc il plonge plus doucement près de la fosse et se redresse en dessous.

\textbf{6.} Les écarts sont $-2.1$, $2.0$, $-7.3$, $5.3$, $7.1$, $2.3$,
$-15.0$ et $8.1$ \unit{km}. Le séisme 7 se trouve à \qty{15}{km} au-dessus de la droite, le séisme 8
à \qty{8}{km} en dessous~; les autres sont à moins d'environ \qty{7}{km} d'elle, une fraction
de l'épaisseur du panneau.

\textbf{7.} $\ang{320.9} - \ang{300} = \ang{20.9}$. Perpendiculairement à la
fosse~: $87.4\cos\ang{20.9} = \qty{81.6}{mm.\year^{-1}}$~; parallèlement~:
$87.4\sin\ang{20.9} = \qty{31.2}{mm.\year^{-1}}$. La convergence est oblique.

\textbf{8.} \qty{81.6}{mm.an^{-1}} fait \qty{81.6}{km.\Ma^{-1}}~: \qty{816}{km} en
\qty{10}{\Ma}.

\textbf{9.} Le chemin est $632/\sin\ang{45} = \qty{894}{km}$, parcouru en
$894/81.6 = \qty{11.0}{\Ma}$.

\textbf{10.} $122 + 11 = \qty{133}{\Ma}$.

\textbf{11.} $t = \qty{10.95}{\Ma} = \qty{3.46e14}{s}$ et
$\sqrt{\num{e-6}\times\num{3.46e14}} = \qty{1.86e4}{m}$, environ \qty{19}{km}
depuis chaque face~: bien moins que les \qty{50}{km} jusqu'au milieu d'un panneau
épais de \qty{100}{km}, dont le cœur est encore froid.

\textbf{12.} Le rapport $q/k$ vaut~: \qty{17.2}{K.km^{-1}} sous le vieux plancher océanique,
\qty{14.0}{K.km^{-1}} sous l'\omterm{def:b1:active-margins:active-margin}{avant-arc} et \qty{32.0}{K.km^{-1}} sous
l'arrière-arc.

\textbf{13.} Sous l'\omterm{def:b1:active-margins:active-margin}{avant-arc} le panneau froid n'est qu'à quelques dizaines de kilomètres
de profondeur~; il absorbe de la chaleur de la roche au-dessus de lui, donc moins de chaleur atteint la
surface.

\textbf{14.} $h = \qty{2.60e-3}{m.s^{-2}}/(2\pi\times\num{6.674e-11}\times1700)
= \qty{3647}{m}$, environ \qty{3.6}{km}~: autant que, ou un peu plus que, toute la
profondeur de la fosse sous le plancher qui l'entoure, environ \qty{3.5}{km}. La
fosse manque d'une masse que rien ne compense.

\textbf{15.} La différence de masse volumique est $3300 - 1000 = \qty{2300}{kg.m^{-3}}$,
et $2\pi\times\num{6.674e-11}\times2300\times3500 = \qty{3.38e-3}{m.s^{-2}}$,
\qty{338}{\milli\gal}. La formule de la plaque compte la masse manquante à toute
distance~; la vraie fosse n'est fléchie vers le bas que sur une bande large de dizaines de kilomètres,
donc son déficit est plus petit, environ \qty{260}{\milli\gal}, entre les deux
estimations de plaque.

\textbf{16.} La résistance de la plaque qui fléchit, qui maintient la fosse vers le bas et
le bombement externe vers le haut comme le ferait une plaque élastique, et la traction du panneau dense
qui plonge dans le \omterm{def:b1:layered-earth:layers}{manteau}.

\textbf{17.} $-61.2 + 1.028\times164 = \qty{107.4}{km}$.

\textbf{18.} Oui~: \qty{107}{km} se situe entre 72 et \qty{173}{km}, et seulement
environ \qty{2}{km} plus profond que la moyenne, \qty{105}{km}.

\textbf{19.} $x = (100 + 61.2)/1.028 = \qty{156.8}{km}$.

\textbf{20.} $\ang{45.8}/2 = \ang{22.9}$, et
$6371\times22.9\times\pi/180 = \qty{2546}{km}$.

\textbf{21.} L'ajustement, \ang{21.9}, est plus petit de \ang{1.0}, environ
\qty{4}{\percent}~: la coque pliée décrit bien la courbure de l'arc des Kouriles.

\textbf{22.} $x = (105 + 61.2)/1.028 = 166.2/1.028 = \qty{161.7}{km}$.

\textbf{23.} $164 - 161.7 = \qty{2.3}{km}$ plus près de la fosse que le front
observé, $2.3/164 = \qty{1.4}{\percent}$~: les volcans des Kouriles se dressent au-dessus d'un panneau
très légèrement plus profond que la moyenne, comme l'a trouvé la question 18.

\textbf{24.} Au kilomètre près, la droite du panneau des Kouriles prédit une
\omterm{def:b1:active-margins:volcanic-arc}{distance arc--fosse} de \textbf{\qty[text-series-to-math=true]{162}{km}} pour
des volcans au-dessus du panneau à \qty{105}{km} (observé~: \qty{164}{km}).
\end{solution}
